\section{Full Metric Definitions and Evaluation Protocols}
\label{app:metric-details}

The 0--3000 keV protocol below applies to the reevaluated MJD and EXO-200
classic models, nine EXO-200 Transformer variants, and six MJD Transformer
variants with MLP or Fourier position encoding. The NEXT and SuperNEMO
classic results and the separate SuperNEMO energy-only illustration retain
their preceding evaluation profile, specified below. The three MJD RoPE
variants and the NEXT and SuperNEMO Transformer entries retain their
historical reported values outside this reevaluation.

\subsection{Energy-matched AUC}
\label{sec:matched-auc}
For event $x_i$, let $y_i\in\{0,1\}$ denote its evaluation label, $E_i$ its
conditioning energy, $s_i$ its scalar classifier score, and $a_i\geq0$ its
base event weight. Scores are oriented so that larger values favor $y_i=1$;
this convention need not coincide with the stored training label
(Table~\ref{tab:metric-protocols}). The selected evaluations use unit base
weights. We first describe the weighted form implemented by the evaluators.
Only finite scores with positive base weight contribute to inclusive AUC;
energy-dependent metrics additionally require finite physical energy satisfying
$0\leq E_i\leq3000$ keV. This range restriction precedes support estimation.

The common support $C=[L,U]$ is the intersection of the two classes'
central energy ranges:
\begin{equation}
 L=\max_{c\in\{0,1\}} Q_c(\alpha),\qquad
 U=\min_{c\in\{0,1\}} Q_c(1-\alpha),\qquad \alpha=0.005,
 \label{eq:common-support}
\end{equation}
where $Q_c$ is the base-weighted empirical energy quantile of finite-score,
finite-energy, positive-weight events in class $c$ within 0--3000 keV.
No additional matching region of interest is specified in these runs.
For energy bins $B_k$, define the count $n_{ck}$ and base mass $m_{ck}$ of
class-$c$ events in $C\cap B_k$, and let
$p_{ck}=m_{ck}/\sum_\ell m_{c\ell}$. A bin is retained only if it contains
at least 20 positive-weight events from each class. Writing $\mathcal V$ for
these bins, the implemented overlap target and event weights are
\begin{align}
 q_k&=\frac{\mathrm{1}\{k\in\mathcal V\}\min(p_{0k},p_{1k})}
 {\sum_{\ell\in\mathcal V}\min(p_{0\ell},p_{1\ell})},
 \label{eq:overlap-target}\\
 w_i&=a_i\frac{q_k}{m_{y_i k}}
 \quad\text{if }E_i\in C\cap B_k,\ k\in\mathcal V;
 \qquad w_i=0\text{ otherwise}.
 \label{eq:matching-weights}
\end{align}
Thus each class has unit total matched weight and assigns mass $q_k$ to
bin $k$. The target favors energy regions represented in both classes and is
defined by the minimum of their class-normalized bin masses.

All retained events are pooled into one weighted ROC curve. Equivalently,
with $H(u)=\mathrm{1}\{u>0\}+\tfrac12\mathrm{1}\{u=0\}$,
\begin{equation}
 \operatorname{AUC}_{\mathrm{match}}=
 \frac{\displaystyle\sum_{i:y_i=1}\sum_{j:y_j=0}
 w_iw_j H(s_i-s_j)}
 {\displaystyle\left(\sum_{i:y_i=1}w_i\right)
 \left(\sum_{j:y_j=0}w_j\right)}.
 \label{eq:matched-auc}
\end{equation}
This is the weighted pairwise ranking interpretation of AUC
\citep{hanley1982meaning}; equal scores receive half credit. The code implements the equivalent weighted pairwise calculation by
sorting scores and processing ties together. It does not average within-bin AUCs: pairs from different
energy bins also contribute to Equation~\ref{eq:matched-auc}.
Inclusive AUC uses the same ranking calculation with base weights and without
energy restriction.

Matching is a deterministic evaluation operation on the supplied split;
it does not resample events or alter the fitted classifier. For the reported
test results, support and weights are estimated on the test events themselves.
The formal matched AUC is withheld if either class retains less than 50\% of
its original finite-score/finite-energy class base mass, or fewer than two
bins survive. Exactly constant common energy exempts only the two-bin
requirement; the per-class event-count and coverage gates still apply.
Empty, single-class, and insufficient-count bins receive zero matched weight.

Both energy metrics in the MJD and EXO-200 reevaluation use the same global
grid after conversion to keV:
\begin{align}
 B_k&=[5k,5(k+1))\ \mathrm{keV},\quad k=0,\ldots,598,\\
 B_{599}&=[2995,3000]\ \mathrm{keV}.
 \label{eq:fixed-bins}
\end{align}
There are exactly 600 energy bins. Exactly 3000 keV belongs to the last bin;
energies below zero or above 3000 keV are excluded and reported separately.
Energies are never clipped. Matching additionally requires
$L\leq E_i\leq U$, so support boundaries may cut through an energy bin.

Matching equalizes bin masses, not continuous energy distributions within
bins. Finite width, finite statistics, trimming, and sparse-bin exclusion
limit the interpretation to the retained support and resolution. A score
that is solely a function of the matched bin index has AUC $1/2$, because
its two weighted class distributions are identical; a score depending on
within-bin energy need not. This is a property of the weighting construction,
not an empirical claim about the trained models.

\subsection{Energy independence score}
\label{sec:energy-independence}
Energy matching measures discrimination after balancing class spectra.
The complementary independence score asks whether the \emph{distribution
of classifier scores} changes with energy within a true class. Conditioning
on class avoids interpreting class separation itself as unwanted dependence.
The evaluator groups by supplied physical category, falling back to the binary
label when category metadata are absent. Here each reported task has two
such groups. This calculation uses the original base weights, without
matching or common-support trimming.

For group $g$, let $P_{gk}$ be the normalized score histogram in an eligible
energy bin, $P_g$ the group's pooled score histogram, and $\pi_{gk}$ the
fraction of retained base weight in that energy bin. Score bins are shared
within a group and obtained from up to 20 weighted score-quantile bins.
Using natural logarithms, define
\begin{align}
 \operatorname{JS}(P,Q)
 &=\tfrac12\sum_r P_r\log\frac{P_r}{M_r}
   +\tfrac12\sum_r Q_r\log\frac{Q_r}{M_r},
 \qquad M_r=\tfrac12(P_r+Q_r),
 \label{eq:jsd}\\
 D_g&=\sum_k\pi_{gk}\operatorname{JS}(P_{gk},P_g),\qquad
 I_g=\operatorname{clip}_{[0,1]}
       \left(1-\sqrt{D_g/\log 2}\right),
 \label{eq:independence-group}\\
 I&=\frac{\sum_{g\in\mathcal G_*} A_g I_g}
           {\sum_{g\in\mathcal G_*} A_g}.
 \label{eq:independence-total}
\end{align}
Here $r$ indexes score-histogram bins, zero-probability terms contribute
zero, $\mathcal G_*$ contains estimable groups, and $A_g$ is their base
weight after finite-pair and 0--3000 keV range filtering, but before
sparse-bin exclusion. These MJD and EXO-200 evaluations
use the retained-event pooling and eligibility rules specified below.
The Jensen--Shannon divergence is bounded by $\log 2$
\citep{lin1991divergence}, giving $I_g,I\in[0,1]$ whenever estimable.
Larger values indicate less detected dependence at the chosen binning
resolution. The report also retains $\min_g I_g$ to expose a poorly behaving
group that the weighted mean could obscure.

This is a histogram-based dependence proxy, not Pearson correlation,
distance correlation, mutual information, or a test certifying statistical
independence. Its ideal value is one when all compared histograms coincide.
Constant predictions also attain one while providing AUC $1/2$; a single
populated energy bin makes the score uninformative about energy dependence.
Finite samples can produce nonzero histogram divergences even under population
independence, and coarse bins can hide dependence. The score depends on energy
coverage, score discretization, and group composition, so it should be compared
under a common protocol. No threshold or operating-point diagnostic enters
Equations~\ref{eq:jsd}--\ref{eq:independence-total}.

\paragraph{Joint interpretation.}
High matched AUC indicates ranking ability after the specified spectral
balancing; high $I$ indicates stable within-class score distributions across
the inspected energy bins. Neither implies the other. Inclusive AUC supplies
the complementary performance under the original sampled spectra. Its change
after matching is descriptive, rather than a causal decomposition of energy
use: both the support and the event weights can change. We report the three
quantities together.

\subsection{Implementation Details}\label{app:detail}
\begin{table}[t]
\caption{Conditioning variables for classic models and the reevaluated EXO-200 and MJD Transformers. The MJD and EXO-200 reevaluation uses 600 fixed 5 keV bins on 0--3000 keV and excludes energies outside this range (v3). NEXT and SuperNEMO classic results retain 600 regular bins plus one overflow bin (v2). The two metrics retain distinct statistical definitions.}
\label{tab:metric-protocols}
\centering\small
\setlength{\tabcolsep}{4pt}
\begin{tabular}{p{.14\linewidth}p{.31\linewidth}p{.45\linewidth}}
\toprule
Dataset & Evaluation-positive class & Physical conditioning energy \\
\midrule
NEXT & $0\nu\beta\beta$ & Summed deposited-hit energy; MeV converted to keV \\
MJD & Clean pulse & Original calibrated float64 energy (keV), recovered from test-event metadata \\
EXO-200 & Stored label 0 ($n_{\mathrm{CCL}}=1$) & Reconstructed rotated energy (keV); negate background-oriented logit \\
SuperNEMO & $2\nu\beta\beta$ & Reconstructed calorimeter $E_1+E_2$ (keV) \\
\bottomrule
\end{tabular}
\end{table}


\paragraph{label selection} The stored NEXT labels assign 1 to $0\nu\beta\beta$ and 0 to
$^{214}\mathrm{Bi}$, while SuperNEMO assigns 1 to $2\nu\beta\beta$ and 0 to
$^{214}\mathrm{Bi}$. EXO-200 stores label 1 for events with more than one
charge cluster and label 0 for single-cluster events; during evaluation, we
reverse the label and score orientation so that single-cluster events form
the positive class, following Table~\ref{tab:metric-protocols}. MJD assigns
label 1 only when all four reference pulse-shape-discrimination criteria
pass. The evaluation-positive class therefore has a dataset-specific
physical interpretation. NEXT uses an approximately $80/10/10$
training/validation/test partition based on class-wise event counts, and
different portions of a source file may enter different splits. SuperNEMO
uses the same approximate proportions while assigning blocks of complete
events within each source category. EXO-200 keeps acquisition runs disjoint
across splits. MJD preserves its official test files and divides the provided
training data $90/10$ into training and validation subsets. Each model for a
given dataset reuses the same recorded partition.

\paragraph{Evaluation inputs and orientation.}
The MJD and EXO-200 reevaluation reuses saved test predictions without
training or checkpoint substitution. The unchanged NEXT classic results
retain stored scalar classifier scores and convert summed deposited-hit
energy from MeV to keV. MJD uses calibrated physical
energy recovered in official test-shard/row order. Event identities and the
saved labels and energies are verified against the original metadata. The
classic adapter's clipping is removed; Transformer energy caches are restored
from float32 to the original float64 values. The three classic binary
classifiers and all six reevaluated Transformers retain their single
clean-oriented scores, while classic GINE retains the recorded aggregation of four PSD outputs
(Appendix~\ref{app:wing-models}). EXO-200's stored label 0 is positive
($n_{\mathrm{CCL}}=1$): we set $y=1-y_{\rm stored}$ and negate the stored
background-oriented logit for both metrics, including all nine reevaluated
Transformers. The unchanged SuperNEMO classic logits favor
$2\nu\beta\beta$ and use reconstructed calorimeter $E_1+E_2$ in keV.
The separate $0\nu\beta\beta$ illustration uses $s=E_1+E_2$ directly.
All selected files have unit base weights. No sigmoid is added to scalar
logits: reproducing $I$ requires retaining the score representation as well
as the ranking direction.

\paragraph{Exact weighted quantiles.}
The evaluator sorts finite observations $z_{(j)}$ with positive weights
$b_{(j)}$ and assigns cumulative center positions
\begin{equation}
 t_j=\frac{\sum_{h\leq j}b_{(h)}-b_{(j)}/2}{\sum_h b_{(h)}}.
\end{equation}
The empirical quantile $Q(u)$ is linearly interpolated between
$(t_j,z_{(j)})$, with endpoint values used outside $[t_1,t_n]$.
Stable sorting preserves input order for tied observations. These are the
quantiles in Equation~\ref{eq:common-support} and in the histogram
construction below. Setting the trim fraction to zero would use each class's
minimum and maximum; the reported configurations instead use 0.005.

\paragraph{Units, boundaries, and population.}
Physical energy is converted to float64 keV before bin assignment. To avoid
a unit-dependent one-step floating-point shift at a theoretical edge,
an input MeV value exactly equal to a canonical keV edge divided by 1000
maps to that exact edge. Other nearby observations are not snapped by a
tolerance. Internal boundaries enter the bin starting at that boundary,
with the explicit 3000 keV exception in Equation~\ref{eq:fixed-bins}.
Original energies are retained in the audit, including MJD's 67 events above
3000 keV (66 nonclean and one clean), excluded from the two energy metrics
rather than clipped into the 2995--3000 keV bin. Inclusive AUC always uses the original finite-score,
positive-weight population, regardless of energy.

\paragraph{Matching calculation and diagnostics.}
Reproduction first retains finite score/energy pairs with
$0\leq E_i\leq3000$ keV and positive base weight, then estimates $C$ and
assigns the shared 600 bins. Event counts and base masses are computed within
$C\cap B_k$. The remaining steps reject bins with fewer than 20 events of
either class, compute
Equations~\ref{eq:overlap-target}--\ref{eq:matching-weights}, and evaluate
one pooled weighted ROC. The class fractions $p_{ck}$ are normalized over
all common-support bins \emph{before} invalid bins are removed. Sparse bins
are not merged. The formal result is unavailable if support is disjoint,
statistics are inadequate, or a reporting gate fails. Coverage for class $c$
is
\begin{equation}
 \gamma_c=\frac{\sum_{i:y_i=c,\ w_i>0}a_i}
 {\sum_{i:y_i=c,\ s_i,E_i\ \mathrm{finite},\ a_i>0}a_i}.
\end{equation}
This is retained \emph{base} mass, not the matched mass, which is normalized
to one in each class. Its denominator precedes both lower and upper energy
range exclusions; range filtering cannot inflate the reported coverage.
Range loss, subsequent support loss and sparse-bin loss are recorded
separately. Both $\gamma_c\geq0.5$ and at least two valid bins are
required. Exactly constant common energy exempts only the two-bin requirement;
the coverage and per-class event-count gates still apply. A diagnostic weighted
AUC may be saved even when the formal reporting gate fails; it is not
substituted for a missing formal result in our table.
The effective sample size for any weights $v_i$ is
$(\sum_i v_i)^2/\sum_i v_i^2$.
Table~\ref{tab:matching-diagnostics} reports support and coverage for the
reference CNN prediction sets.

\begin{table}[!htbp]
\centering\footnotesize
\caption{Matching diagnostics for the multi-view CNN prediction sets under their dataset-specific profiles (Table~\ref{tab:metric-protocols}). All reevaluated classic models within each displayed test population use identical test event IDs, labels, physical energies, and base weights. Coverage $\gamma_c$ uses the original finite score/energy class base mass before energy-range exclusion; ESS uses matched weights.}
\label{tab:matching-diagnostics}
\setlength{\tabcolsep}{4pt}
\resizebox{\linewidth}{!}{%
\begin{tabular}{lrrrrrr}
\toprule
Dataset & Common support (keV) & Bins & $\gamma_1$ & $\gamma_0$ & ESS$_1$ & ESS$_0$ \\
\midrule
NEXT & [2434.84, 2480.35] & 11 & 0.9882 & 0.9700 & 55195 & 42260 \\
MJD & [18.15, 2613.79] & 493 & 0.9874 & 0.9631 & 141205 & 213947 \\
EXO-200 & [538.18, 2662.57] & 426 & 0.9457 & 0.9414 & 51924 & 61627 \\
SuperNEMO & [350.21, 2218.25] & 374 & 0.9900 & 0.8988 & 312506 & 207113 \\
\bottomrule
\end{tabular}}
\end{table}


\begin{table}[!htbp]
\centering\footnotesize
\caption{Disjoint matching losses and the separate independence sparse-bin loss, reported as positive/negative event counts. $N_1/N_0$ is the original test population. Range loss precedes support estimation, then matching sparse-bin exclusion. For MJD and EXO-200, range loss includes energies below 0 or above 3000 keV; the $E>3000$ column is a subset of that loss. For retained NEXT and SuperNEMO results, only negative energies are excluded and $E>3000$ counts are retained overflow diagnostics. The separate $I$ sparse loss is measured after the applicable range filtering. Counts are shared by reevaluated classic models within each displayed test population.}
\label{tab:evaluation-losses}
\setlength{\tabcolsep}{4pt}
\resizebox{\linewidth}{!}{%
\begin{tabular}{lrrrrrr}
\toprule
Dataset & $N_1/N_0$ & Range loss & $E>3000$ & $I$ sparse loss & Support loss & Matching sparse loss \\
\midrule
NEXT & 64879/51670 & 0/0 & 0/0 & 0/0 & 767/1550 & 0/0 \\
MJD & 148251/241749 & 1/66 & 1/66 & 137/214 & 1746/7378 & 121/1481 \\
EXO-200 & 66062/74321 & 3/1423 & 3/1423 & 152/257 & 3587/2932 & 0/0 \\
SuperNEMO & 327680/262144 & 0/0 & 0/40 & 333/271 & 3275/26530 & 0/0 \\
\bottomrule
\end{tabular}}
\end{table}


No resampling, new model fitting, bootstrap interval, or multiple-seed average
is introduced. Support restriction and matching weights can both alter the
AUC relative to the inclusive population; their difference is descriptive,
not a causal attribution to energy.

\paragraph{Fixed-grid independence score.}
The reevaluated MJD and EXO-200 records group events by physical category when supplied, otherwise by
binary label. Within each group the evaluator removes nonfinite
score/energy pairs, energies outside 0--3000 keV, and zero-weight events, retains only
energy bins with at
least 20 events, and computes weighted score quantiles using the retained
events. Duplicate score edges are removed. For nonconstant scores the extreme
edges are moved outward by one floating-point step; constant scores receive
one interval centered at that value with half-width
$\max(|s|10^{-9},10^{-12})$. Score histograms use left-closed/right-open bins
with the last right edge included. They are normalized probability masses,
not densities. No pseudocount is added.

Let $R_g$ be the retained events, $R_{gk}$ those in bin $k$, and $S_{gr}$
score-histogram interval $r$. Then
\begin{align}
 P_{gk,r}&=\frac{\sum_{i\in R_{gk}}a_i\mathrm{1}\{s_i\in S_{gr}\}}
 {\sum_{i\in R_{gk}}a_i},\qquad
 P_{g,r}=\frac{\sum_{i\in R_g}a_i\mathrm{1}\{s_i\in S_{gr}\}}
 {\sum_{i\in R_g}a_i},\label{eq:score-histograms}\\
 \pi_{gk}&=\frac{\sum_{i\in R_{gk}}a_i}{\sum_{i\in R_g}a_i}.
\end{align}
Equations~\ref{eq:jsd}--\ref{eq:independence-group} operate on these
histograms. The final group weight $A_g$ includes all finite-pair, positive-weight
events in that group within 0--3000 keV, including events from rejected sparse bins. Thus
the overall score is not an unweighted class mean and not a mean weighted
only by retained events. A group with fewer than 20 finite events or no
eligible energy bin is excluded; if no group is estimable, the score is
missing. One eligible bin is sufficient for this implementation, unlike
the usual two-bin reporting gate for matched AUC.

\begin{table}[!htbp]
\centering\footnotesize
\caption{Class-level independence diagnostics under the profiles in Table~\ref{tab:metric-protocols}. Retained fractions divide eligible-bin base mass by original finite score/energy class base mass, before energy-range exclusion. Score histograms retain up to 20 weighted score-quantile bins. The audit files include all extended NEXT configurations and full precision.}
\label{tab:independence-diagnostics}
\setlength{\tabcolsep}{4pt}
\resizebox{\linewidth}{!}{%
\begin{tabular}{lrrrrr}
\toprule
Dataset / model & $I_1$ & $I_0$ & $\min_g I_g$ & Bins$_{1/0}$ & Retained$_{1/0}$ \\
\midrule
NEXT / Multi-view CNN & 0.9763 & 0.9643 & 0.9643 & 11/11 & 1.0000/1.0000 \\
NEXT / Static GINE & 0.9778 & 0.9609 & 0.9609 & 11/11 & 1.0000/1.0000 \\
NEXT / BiGRU & 0.9753 & 0.9657 & 0.9657 & 11/11 & 1.0000/1.0000 \\
NEXT / PointMamba-lite & 0.9776 & 0.9632 & 0.9632 & 11/11 & 1.0000/1.0000 \\
MJD / Multi-view CNN & 0.7325 & 0.7763 & 0.7325 & 496/524 & 0.9991/0.9988 \\
MJD / Static GINE & 0.8057 & 0.7911 & 0.7911 & 496/524 & 0.9991/0.9988 \\
MJD / BiGRU & 0.7791 & 0.7730 & 0.7730 & 496/524 & 0.9991/0.9988 \\
MJD / PointMamba-lite & 0.7519 & 0.7865 & 0.7519 & 496/524 & 0.9991/0.9988 \\
EXO-200 / Multi-view CNN & 0.7315 & 0.7709 & 0.7315 & 438/484 & 0.9977/0.9774 \\
EXO-200 / Static GINE & 0.7925 & 0.8093 & 0.7925 & 438/484 & 0.9977/0.9774 \\
EXO-200 / BiGRU & 0.7739 & 0.8033 & 0.7739 & 438/484 & 0.9977/0.9774 \\
EXO-200 / PointMamba-lite & 0.7854 & 0.7863 & 0.7854 & 438/484 & 0.9977/0.9774 \\
SuperNEMO / Multi-view CNN & 0.9213 & 0.8962 & 0.8962 & 417/514 & 0.9990/0.9990 \\
SuperNEMO / Static GINE & 0.9241 & 0.8981 & 0.8981 & 417/514 & 0.9990/0.9990 \\
SuperNEMO / BiGRU & 0.9292 & 0.9105 & 0.9105 & 417/514 & 0.9990/0.9990 \\
SuperNEMO / PointMamba-lite & 0.9292 & 0.9119 & 0.9119 & 417/514 & 0.9990/0.9990 \\
\bottomrule
\end{tabular}}
\end{table}


\begin{table}[!htbp]
\centering\footnotesize
\caption{Diagnostics for the 15 reevaluated MJD and EXO-200 Transformers using 600 fixed 5 keV bins on 0--3000 keV. All match the CNN test identities, class partitions, physical energies and unit weights within their dataset, so common support and matching losses equal Tables~\ref{tab:matching-diagnostics}--\ref{tab:evaluation-losses}. Retained fractions and coverage use original finite-pair class mass before range exclusion. The minimum group score is the smaller of the two displayed $I_g$ values. All formal matched AUCs pass the reporting gates.}
\label{tab:transformer-diagnostics}
\setlength{\tabcolsep}{4pt}
\resizebox{\linewidth}{!}{%
\begin{tabular}{lrrrrrr}
\toprule
Dataset / Transformer & $I_1/I_0$ & Bins$_{I,1/0}$ & Retained$_{I,1/0}$ & $K_{\rm match}$ & $\gamma_1/\gamma_0$ & ESS$_{1/0}$ \\
\midrule
EXO-200 / entity--mlp & 0.6868/0.7828 & 438/484 & 0.9977/0.9774 & 426 & 0.9457/0.9414 & 51924/61627 \\
EXO-200 / entity--fourier & 0.6545/0.7969 & 438/484 & 0.9977/0.9774 & 426 & 0.9457/0.9414 & 51924/61627 \\
EXO-200 / entity--rope & 0.6583/0.7671 & 438/484 & 0.9977/0.9774 & 426 & 0.9457/0.9414 & 51924/61627 \\
EXO-200 / region--mlp & 0.7893/0.8115 & 438/484 & 0.9977/0.9774 & 426 & 0.9457/0.9414 & 51924/61627 \\
EXO-200 / region--fourier & 0.7646/0.8131 & 438/484 & 0.9977/0.9774 & 426 & 0.9457/0.9414 & 51924/61627 \\
EXO-200 / region--rope & 0.7502/0.8246 & 438/484 & 0.9977/0.9774 & 426 & 0.9457/0.9414 & 51924/61627 \\
EXO-200 / summary--mlp & 0.7296/0.8183 & 438/484 & 0.9977/0.9774 & 426 & 0.9457/0.9414 & 51924/61627 \\
EXO-200 / summary--fourier & 0.7894/0.8132 & 438/484 & 0.9977/0.9774 & 426 & 0.9457/0.9414 & 51924/61627 \\
EXO-200 / summary--rope & 0.7713/0.7740 & 438/484 & 0.9977/0.9774 & 426 & 0.9457/0.9414 & 51924/61627 \\
MJD / entity--mlp & 0.8388/0.7950 & 496/524 & 0.9991/0.9988 & 493 & 0.9874/0.9631 & 141205/213947 \\
MJD / entity--fourier & 0.8196/0.7913 & 496/524 & 0.9991/0.9988 & 493 & 0.9874/0.9631 & 141205/213947 \\
MJD / region--mlp & 0.7935/0.8013 & 496/524 & 0.9991/0.9988 & 493 & 0.9874/0.9631 & 141205/213947 \\
MJD / region--fourier & 0.7869/0.7881 & 496/524 & 0.9991/0.9988 & 493 & 0.9874/0.9631 & 141205/213947 \\
MJD / summary--mlp & 0.8019/0.7973 & 496/524 & 0.9991/0.9988 & 493 & 0.9874/0.9631 & 141205/213947 \\
MJD / summary--fourier & 0.7914/0.7759 & 496/524 & 0.9991/0.9988 & 493 & 0.9874/0.9631 & 141205/213947 \\
\bottomrule
\end{tabular}}
\end{table}


\paragraph{Changes from historical implementations.}
The historical EXO-200 evaluations, including the nine Transformer variants,
used six class-balanced
pooled energy-quantile bins for matching, and a separate eight per-category
energy-quantile construction for $I$. Both energy constructions are replaced
by Equation~\ref{eq:fixed-bins}; the up-to-20 score-quantile bins remain.
Their old $I$ implementation could compare a local energy bin containing
only two events and pooled all finite events. The common reference instead
requires 20 events per group/bin and computes score edges and the pooled
histogram only on retained events. Bin weights use retained group mass;
final group weights use eligible population mass before sparse exclusion.
These choices follow the former NEXT/MJD fixed-grid estimator and preserve
the JS-then-square-root and final group-aggregation definitions.
The six MJD Transformer caches used 836 fixed 5 keV bins extending to
4180 keV, without clipping. They now use Equation~\ref{eq:fixed-bins} and
the same range exclusions as the other reevaluated models. Their original
float64 physical energies are restored before binning. These changes can
alter both score-histogram populations and support quantiles, even though
the historical MJD Transformer bin width was already 5 keV.

\paragraph{Retained NEXT and SuperNEMO profile.}
The previously reported NEXT and SuperNEMO classic results, including NEXT's
separate exploratory batches, and the SuperNEMO energy-only illustration
retain \texttt{EnergyBench-unified-}\allowbreak\texttt{5keV-overflow-v2.0.0}.
This profile uses the 600 regular bins in Equation~\ref{eq:fixed-bins} plus
$B_{600}=(3000,\infty)$ keV. Negative energies are excluded, while all
finite nonnegative energies enter the support-quantile calculation and the
pre-sparse group masses $A_g$. The overflow bin is subject to the same
event-count rules as a regular bin. The score-quantile histograms, overlap
target, weighted AUC and original finite-pair coverage denominators use
the same definitions as above. These retained results and their figures
are not recomputed under the MJD/EXO-200 range restriction.

\paragraph{Reproduction and validation.}
The MJD/EXO-200 configuration is \texttt{EnergyBench-unified-}\allowbreak\texttt{5keV-range-v3.0.0}.
The preceding v2 configuration remains available for the retained NEXT and
SuperNEMO results; each result records its applicable profile and SHA256
fingerprint. The reproducibility bundle contains the corresponding frozen
evaluators, event-input manifests and hashes, full-precision
metrics, class-level bin counts and losses, and a table generator that reads
these results. Numerical checks cover unit equivalence at all 601 finite
edges, adjacent floating-point values, exclusion outside 0--3000 keV, constants, ties,
missing and sparse inputs, weighted pair enumeration, and replay of the
historical fixed-grid estimator on identical supported inputs. Independent
recalculation checks the metrics, coverage gates, and original event sources.

\paragraph{Degeneracies and scope.}
The unified classification entry point rejects labels outside $\{0,1\}$
on the entire input array. It also rejects negative or nonfinite weights and
misaligned arrays; entirely zero weight yields unestimable metrics. Nonfinite scores are excluded
from ROC, and missing energy is additionally excluded from energy-dependent
statistics. An absent ROC class makes AUC unavailable; $I$ can still be
computed for estimable groups. If no category array is supplied, binary labels
define the groups; missing entries in a supplied category array are rejected.
The dependence routine filters nonfinite pairs; no correlation coefficient
enters $I$. Constant scores produce identical histograms
and $I=1$. Constant conditioning energy likewise supplies no meaningful
test of energy dependence.

The score in Equation~\ref{eq:independence-group} averages pairwise
Jensen--Shannon divergences \emph{before} taking a square root; it is not an
average of Jensen--Shannon distances and not the multi-distribution
Jensen--Shannon divergence. We use a descriptive benchmark name for this
implemented transformation without claiming that it is a standard test or
an original statistical construction.
